\begin{answer}
The classical (unsupervised) EM algorithm is implemented in \texttt{run\_em} in \texttt{src/semi\_supervised\_em/gmm.py}. The algorithm is initialized with random group-based means and covariances, and iterates E- and M-steps until the log-likelihood change falls below $\varepsilon = 10^{-3}$.

Three scatter plots (one per trial) are generated by the \texttt{plot\_gmm\_preds} function. Each point is colored by its most probable cluster assignment (argmax of the final $w$ matrix). Due to random initialization, the cluster color assignments may vary across trials, but the geometric cluster structure should be roughly consistent when initialization is favorable.

\textbf{Observation:} The unsupervised EM is sensitive to initialization. With some random seeds, two or more clusters may collapse onto the same Gaussian, resulting in poor separation, especially for the high-variance Gaussian (cluster 4). This is expected because without labeled data the algorithm cannot distinguish closely overlapping clusters reliably.

\textit{[See the generated PDF plots \texttt{pred\_0.pdf}, \texttt{pred\_1.pdf}, \texttt{pred\_2.pdf}.]}
\end{answer}
